% Part: first-order-logic
% Chapter: sequent-calculus
% Section: derivations

\documentclass[../../../include/open-logic-section]{subfiles}

\begin{document}

\iftag{FOL}
      {\olfileid[fr]{fol}{seq}{der}}
      {\olfileid[fr]{pl}{seq}{der}}

\olsection{\usetoken{P}{derivation}}

\begin{explain}
Nous avons précisé la forme des séquents initiaux et donné les règles
d'inférence. Les !!{derivation}s du calcul des séquents sont
engendrées inductivement à partir de ces éléments : toute
!!{derivation} est soit un séquent initial à lui seul, soit formée
d'une ou de deux !!{derivation}s suivies d'une inférence.
\end{explain}

\begin{defn}[\usetoken{S}{derivation} dans $\Log{LK}$]
Une \emph{!!{derivation} dans $\Log{LK}$} d'un séquent~$S$ est un
arbre fini de séquents qui satisfait aux conditions suivantes :
\begin{enumerate}
\item Les séquents situés tout en haut de l'arbre sont des séquents initiaux.
\item Le séquent situé tout en bas de l'arbre est~$S$.
\item Tout séquent de l'arbre autre que $S$ est une prémisse d'une
  application correcte d'une règle d'inférence dont la conclusion
  se trouve immédiatement sous ce séquent dans l'arbre.
\end{enumerate}
On dit alors que $S$ est le \emph{séquent final} de la
!!{derivation} et que $S$ est \emph{!!{derivable} dans $\Log{LK}$}
(ou $\Log{LK}$-!!{derivable}).
\end{defn}

\begin{ex}
Tout séquent initial, par exemple $!C \Sequent !C$, est
!!a{derivation}. On peut en obtenir une nouvelle !!{derivation}
en appliquant, par exemple, la règle $\LeftR{\Weakening}$,
\begin{prooftree}
\Axiom$ \Gamma \fCenter \Delta $
\RightLabel{\LeftR{\Weakening}}
\UnaryInf$ !A, \Gamma \fCenter \Delta$
\end{prooftree}
Cette règle est toutefois générale : on peut y remplacer $!A$ par
n'importe quel !!{sentence}, par exemple par~$!D$. Si la prémisse
coïncide avec notre séquent initial $!C \Sequent !C$, alors
$\Gamma$ et $\Delta$ se réduisent tous deux à~$!C$, et la
conclusion est $!D, !C \Sequent !C$. L'arbre suivant est donc
!!a{derivation} :
\begin{prooftree}
\Axiom$ !C \fCenter !C $
\RightLabel{\LeftR{\Weakening}}
\UnaryInf$ !D, !C \fCenter !C$
\end{prooftree}
On peut maintenant appliquer une autre règle, par exemple
$\LeftR{\Exchange}$, qui permet d'intervertir deux !!{sentence}s
à gauche. L'arbre suivant est donc lui aussi une !!{derivation}
correcte :
\begin{prooftree}
\Axiom$ !C \fCenter !C $
\RightLabel{\LeftR{\Weakening}}
\UnaryInf$ !D, !C \fCenter !C$
\RightLabel{\LeftR{\Exchange}}
\UnaryInf$ !C, !D \fCenter !C$
\end{prooftree}
Dans cette application de la règle, présentée sous la forme
\begin{prooftree}
\Axiom$ \Gamma, !A, !B, \Pi \fCenter \Delta $
\RightLabel{\LeftR{\Exchange}}
\UnaryInf$ \Gamma, !B, !A, \Pi \fCenter \Delta,$
\end{prooftree}
$\Gamma$ et $\Pi$ sont tous deux vides, $\Delta$ est $!C$, et
les rôles de $!A$ et de $!B$ sont joués respectivement par $!D$
et~$!C$. De façon analogue, on voit aussi que
\begin{prooftree}
\Axiom$ !D \fCenter !D $
\RightLabel{\LeftR{\Weakening}}
\UnaryInf$ !C, !D \fCenter !D$
\end{prooftree}
est !!a{derivation}. On peut à présent combiner ces deux
!!{derivation}s au moyen de $\RightR{\land}$. Rappelons cette règle :
\begin{prooftree}
\Axiom$\Gamma \fCenter \Delta, !A$
\Axiom$ \Gamma \fCenter \Delta, !B$
\RightLabel{\RightR{\land}}
\BinaryInf$ \Gamma \fCenter \Delta, !A \land !B $
\end{prooftree}
Dans notre cas, les prémisses doivent coïncider avec les derniers
séquents des !!{derivation}s qui s'y terminent. Cela signifie que
$\Gamma$ est $!C, !D$, que $\Delta$ est vide, que $!A$ est $!C$
et que $!B$ est $!D$. Pour que l'inférence soit correcte, sa
conclusion doit donc être $!C, !D \Sequent !C \land !D$.
\begin{prooftree}
\Axiom$ !C \fCenter !C $
\RightLabel{\LeftR{\Weakening}}
\UnaryInf$ !D, !C \fCenter !C$
\RightLabel{\LeftR{\Exchange}}
\UnaryInf$ !C, !D \fCenter !C$
\Axiom$ !D \fCenter !D $
\RightLabel{\LeftR{\Weakening}}
\UnaryInf$ !C, !D \fCenter !D$
\RightLabel{\RightR{\land}}
\BinaryInf$ !C, !D \fCenter !C \land !D $
\end{prooftree}
On peut bien sûr intervertir les prémisses : $!A$ serait alors
$!D$ et $!B$ serait~$!C$.
\begin{prooftree}
\Axiom$ !D \fCenter !D $
\RightLabel{\LeftR{\Weakening}}
\UnaryInf$ !C, !D \fCenter !D$
\Axiom$ !C \fCenter !C $
\RightLabel{\LeftR{\Weakening}}
\UnaryInf$ !D, !C \fCenter !C$
\RightLabel{\LeftR{\Exchange}}
\UnaryInf$ !C, !D \fCenter !C$
\RightLabel{\RightR{\land}}
\BinaryInf$ !C, !D \fCenter !D \land !C $
\end{prooftree}
\end{ex}

\end{document}
